% Image credits for the photographs and AI illustrations of Book 3 (University, Year 1), Arabic edition.
% Machine-readable license records live in images/book3/CREDITS.md; keep this
% file in sync with frontmatter/image-credits-book3.tex (one \item per item there).
% As in the Book 2 Arabic credits: personal names are transliterated, with the
% original spelling in parentheses (the credit itself); institutions, datasets,
% works and licences stay verbatim inside \babelsublr; only the gloss is translated.
\chapter*{مصادر الصور}

الرسوم والخرائط والمخططات في هذا الكتاب أصلية؛ وما رُسم منها من بيانات
منشورة يذكر مصدره في التعليق المرافق أو في هذه القائمة. أما الصور
الفوتوغرافية فمستعملة بمقتضى رخصها الحرة، مع الشكر لأصحابها؛ وروابط المصادر
الكاملة وروابط الرخص لكل صورة مدرجة في \texttt{images/book3/CREDITS.md} في
مستودع الكتاب.

\bigskip
وإلى جانب الصور الفوتوغرافية والرسوم الأصلية، بعض الأشكال رسوم توضيحية
مولَّدة بالذكاء الاصطناعي، أُنتجت بأداة توليد الصور من OpenAI انطلاقًا من
أوامر كتبها محررو الكتاب، ورُوجعت دقتها العلمية قبل إدراجها. وهي لا تُظهر
إلا مشاهد عامة، ولا تمثّل أبدًا مكانًا حقيقيًا أو شخصًا حقيقيًا أو عيّنة
حقيقية. والأمر الكامل لكل صورة مولَّدة مسجَّل في
\texttt{images/book3/ai/PROMPTS.md} في المستودع.

\bigskip
\noindent الصور الفوتوغرافية والبيانات، فصلًا فصلًا:
\begin{itemize}\small
  \item الفصل~1 (الفلك الكروي): مسارات النجوم فوق مرصد \babelsublr{Paranal Observatory}، صورة \babelsublr{R.~Wesson/ESO}، \babelsublr{CC~BY~4.0}.
  \item الفصل~2 (التلسكوبات والكواشف والقياس الضوئي): ست قطع من المرآة الأولية لتلسكوب جيمس ويب الفضائي قبل اختبارات التبريد، \babelsublr{NASA/MSFC/David Higginbotham}، ملك عام.
  \item الفصل~3 (علم الكواكب المقارن): الكواكب الأرضية بمقياس واحد، \babelsublr{NASA}؛ والزهرة من رادار \babelsublr{Magellan}، \babelsublr{NASA/JPL}؛ والمريخ من المدار \babelsublr{Viking 1}، \babelsublr{NASA/USGS}؛ وأوروبا، \babelsublr{NASA/JPL/DLR}؛ وآيو، \babelsublr{NASA/JPL/University of Arizona}؛ كلها ملك عام.
  \item الفصل~4 (الخرائط الجيولوجية): ويليام سميث (\babelsublr{William Smith})، \babelsublr{\emph{A Delineation of the Strata of England and Wales, with part of Scotland}}، 1815؛ ملك عام (عبر \babelsublr{Wikimedia Commons}).
  \item الفصل~5 (المقاطع الجيولوجية): عدم التوافق عند \babelsublr{Siccar Point} الذي وصفه هاتون (\babelsublr{Hutton})، \babelsublr{PawełS}، \babelsublr{CC~BY-SA~4.0} (عبر \babelsublr{Wikimedia Commons}).
  \item الفصل~6 (طرائق الحقل والمختبر I): شريحة رقيقة من الغرانيت البيوتيتي بين قطبين متقاطعين، دارلا ساندرول (\babelsublr{Darla Sondrol})، \babelsublr{CC0} (عبر \babelsublr{Wikimedia Commons}).
  \item الفصل~7 (المعادن المكوِّنة للصخور): عينات من الأوليفين والكوارتز والبيوتيت والألماندين والهورنبلند، روبرت م.~لافينسكي (\babelsublr{Robert M.~Lavinsky})، \babelsublr{CC~BY-SA~3.0}؛ واللابرادوريت، تويلا بيكر (\babelsublr{Twyla Baker})، \babelsublr{CC0}؛ وبلورات معينية من تشقّق الكالسيت، جيوفانا كانو وإيفا سانتيني (\babelsublr{Giovanna Canu and Eva Santini})، \babelsublr{CC~BY~4.0} (كلها عبر \babelsublr{Wikimedia Commons}).
  \item الفصل~8 (علم المعادن البصري): شريحة رقيقة من الغابرو بقطب واحد وبين قطبين متقاطعين، \babelsublr{Chmee2}، \babelsublr{CC~BY-SA~3.0}؛ وأوغست ميشيل-ليفي (\babelsublr{Auguste Michel-Lévy}) طالبًا، ملك عام (عبر \babelsublr{Wikimedia Commons}). ألوان التداخل محسوبة من دوال مطابقة الألوان \babelsublr{CIE 1931} ومن المُضيء القياسي \babelsublr{CIE standard illuminant D65}، \babelsublr{\copyright{} CIE}، \babelsublr{CC~BY-SA~4.0}.
  \item الفصل~9 (الأرض ذات الطبقات): أندريا موهوروفيتش (\babelsublr{Andrija Mohorovičić})، ملك عام؛ وإنغه ليمان (\babelsublr{Inge Lehmann}) عام 1932، صورة لـ\babelsublr{Even Neuhaus}، ملك عام (عبر \babelsublr{Wikimedia Commons}). نموذج \babelsublr{PREM}: دزيفونسكي وأندرسون (\babelsublr{A.~M.~Dziewonski and D.~L.~Anderson}، 1981)، كما وردت في \babelsublr{ObsPy}.
  \item الفصل~10 (ميزانية حرارة الأرض): تدفّق الحرارة العالمي عن ديفيز وديفيز (\babelsublr{J.~H.~Davies and D.~R.~Davies}، 2010)، \babelsublr{\emph{Solid Earth}}، \babelsublr{CC~BY~3.0}؛ والتدرّج الأديباتي للوشاح عن كاتسورا (\babelsublr{T.~Katsura}، 2022)، \babelsublr{\emph{JGR Solid Earth}}، \babelsublr{CC~BY}.
  \item الفصل~11 (التوازن الأيزوستاتي وشذوذات الجاذبية): جهاز البندول لفينينغ مينيس (\babelsublr{Vening Meinesz})، \babelsublr{Broodt}، \babelsublr{CC~BY-SA~4.0} (عبر \babelsublr{Wikimedia Commons}). الارتفاعات من \babelsublr{ETOPO1} (\babelsublr{NOAA National Centers for Environmental Information})، مأخوذة عيّناتها عند درجة واحدة.
  \item الفصل~12 (حركيّة الصفائح على الكرة): محطة دائمة لتحديد المواقع بالأقمار الصناعية، \babelsublr{Alpsdake}، \babelsublr{CC~BY-SA~4.0} (عبر \babelsublr{Wikimedia Commons}). حدود الصفائح من \babelsublr{PB2002} (ب.~بيرد، \babelsublr{P.~Bird}، 2003؛ بصيغة \babelsublr{GeoJSON} من ه.~أهلينيوس، \babelsublr{H.~Ahlenius}، \babelsublr{ODC-By 1.0})؛ والسواحل من \babelsublr{Natural Earth} (ملك عام).
  \item الفصل~13 (الشقوق والحيود والهوامش السلبية): شقوق مفتوحة في ثينغفيلير (\babelsublr{Þingvellir})، أيسلندا، أولغا إرنست (\babelsublr{Olga Ernst})، \babelsublr{CC~BY-SA~4.0}؛ ولابات وسائدية على صدع غالاباغوس، \babelsublr{NOAA Okeanos Explorer Program}، \babelsublr{CC~BY~2.0} (كلاهما عبر \babelsublr{Wikimedia Commons}). الأعماق من \babelsublr{ETOPO1} (\babelsublr{NOAA National Centers for Environmental Information})، وأعمار قاع المحيط من سيتون وآخرين (\babelsublr{M.~Seton and others}، 2020)، \babelsublr{EarthByte}، \babelsublr{CC~BY~4.0}.
  \item الفصل~14 (الهوامش النشطة): ساريتشيف بيك (\babelsublr{Sarychev Peak}) يثور، كما رُئي من محطة الفضاء الدولية، \babelsublr{NASA}، ملك عام (عبر \babelsublr{Wikimedia Commons}). الزلازل من \babelsublr{USGS ANSS Comprehensive Catalog} (ملك عام)؛ والبراكين من \babelsublr{Smithsonian Institution Global Volcanism Program}؛ وحدود الصفائح من \babelsublr{PB2002} (ب.~بيرد، \babelsublr{P.~Bird}، 2003؛ بصيغة \babelsublr{GeoJSON} من ه.~أهلينيوس، \babelsublr{H.~Ahlenius}، \babelsublr{ODC-By 1.0})؛ والسواحل من \babelsublr{Natural Earth} (ملك عام)؛ وأعمار قاع المحيط من سيتون وآخرين (\babelsublr{M.~Seton and others}، 2020)، \babelsublr{EarthByte}، \babelsublr{CC~BY~4.0}.
  \item الفصل~15 (البقع الساخنة والبراكين داخل الصفائح): نافورة لابة في \babelsublr{Pu\textquoteleft u \textquoteleft O\textquoteleft o}، كيلاويا، \babelsublr{G.~E.~Ulrich}، \babelsublr{U.S. Geological Survey}، ملك عام؛ والغات الغربية في ماثيران، \babelsublr{Nicholas (Nichalp)}، \babelsublr{CC~BY-SA~2.5} (كلاهما عبر \babelsublr{Wikimedia Commons}). مواقع البقع الساخنة وسموت آثارها ومعدلاتها من الجدول~1 في الملحق الإلكتروني لـ\babelsublr{W.~J.~Morgan and J.~P.~Morgan} (2007)؛ والأعماق من \babelsublr{ETOPO1} (\babelsublr{NOAA National Centers for Environmental Information})؛ وحدود الصفائح من \babelsublr{PB2002} (ب.~بيرد، \babelsublr{P.~Bird}، 2003؛ بصيغة \babelsublr{GeoJSON} من ه.~أهلينيوس، \babelsublr{H.~Ahlenius}، \babelsublr{ODC-By 1.0})؛ والسواحل من \babelsublr{Natural Earth} (ملك عام).
  \item الفصل~16 (الإجهاد والانفعال والريولوجيا): خلية ثلاثية المحاور عالية الضغط لقياس مقاومة الصخور، داريل بيروني (\babelsublr{Darryl Peroni})، \babelsublr{CSIRO}، \babelsublr{CC~BY~3.0} (عبر \babelsublr{Wikimedia Commons}). معاملات زحف الانخلاع للكوارتزيت الرطب (\babelsublr{Gleason and Tullis}، 1995) والأوليفين الجاف (\babelsublr{Hirth and Kohlstedt}، 2003) كما جمعها م.~أ.~لوبيز-سانشيز (\babelsublr{M.~A.~Lopez-Sanchez}) في \babelsublr{\emph{strength\_envelopes}} (\babelsublr{MPL-2.0}).
  \item الفصل~17 (الصدوع والطيّات والنسيج الصخري): سطح انزلاق عليه خدوش عند \babelsublr{Corona Heights}، سان فرانسيسكو، \babelsublr{Dextersinister}، \babelsublr{CC~BY-SA~4.0}؛ والبريشيا الصدعية لصدع كيستون الزاحف، نيفادا، \babelsublr{Jstuby}، ملك عام؛ وطيّات الشيفرون في \babelsublr{Millook}، كورنوال، \babelsublr{Smalljim}، \babelsublr{CC~BY~3.0}؛ وطيّة محدّبة مع تفلّق محوري مستوٍ قرب بوتِس، قانتابريا، \babelsublr{Mikenorton}، \babelsublr{CC~BY-SA~3.0} (كلها عبر \babelsublr{Wikimedia Commons}).
  \item الفصل~18 (مصدر الزلزال): أرض وأسيجة انفصلت بالصدع عام 1906 في مقاطعة مارين، \babelsublr{\emph{Popular Science Monthly}}، المجلد~69 (1906)، ملك عام؛ ومجرى \babelsublr{Wallace Creek} المزاح عبر صدع سان أندرياس، \babelsublr{R.~E.~Wallace}، \babelsublr{U.S. Geological Survey}، ملك عام (كلاهما عبر \babelsublr{Wikimedia Commons}). الزلازل من \babelsublr{USGS ANSS Comprehensive Catalog} (ملك عام)؛ والآلية البؤرية لزلزال توهوكو عام 2011 من \babelsublr{USGS W-phase moment tensor} (ملك عام).
  \item الفصل~19 (التجوية والطين): صخر متحلل ورمل جرانيتي عند تارغاسون، جبال البيرينيه، \babelsublr{Pinpin}، \babelsublr{CC~BY~2.5}؛ والبوكسيت على حجر رملي كاولينيتي في ويبا، أستراليا، فيرنر شيلمان (\babelsublr{Werner Schellmann})، \babelsublr{CC~BY-SA~2.5}؛ وبلورات الكاولينيت تحت المجهر الإلكتروني الماسح، \babelsublr{U.S. Geological Survey}، ملك عام (كلها عبر \babelsublr{Wikimedia Commons}). أنصاف الأقطار الأيونية عن \babelsublr{R.~D.~Shannon} (1976).
  \item الفصل~20 (التعرية والنقل): انتشار مياه فيضان المسيسيبي من جهاز \babelsublr{MODIS} على الساتل \babelsublr{Terra}، 4 مارس 2018، صورة لجوشوا ستيفنز (\babelsublr{Joshua Stevens})، \babelsublr{NASA Earth Observatory}، ملك عام (عبر \babelsublr{Wikimedia Commons}).
  \item الفصل~21 (الرسوبيات الفتاتية وبيئاتها): تطبّق متقاطع في حجر آزتيك الرملي، نيفادا، شيب كارون (\babelsublr{Chip Carroon})، \babelsublr{Bureau of Land Management}، \babelsublr{CC~BY-SA~2.0}؛ ودلتا نهر لينا من \babelsublr{Landsat~7}، \babelsublr{U.S. Geological Survey}، ملك عام؛ وطبقة عكرة (\babelsublr{turbidite}) تُظهر أقسام تعاقب \babelsublr{Bouma}، بيكي-أوزه (\babelsublr{Becke-Oese})، جو ويبر (\babelsublr{Jo Weber})، \babelsublr{CC~BY-SA~3.0} (كلها عبر \babelsublr{Wikimedia Commons}).
  \item الفصل~22 (الرسوبيات الكربوناتية والسيليسية والتبخّرية): بورا بورا من محطة الفضاء الدولية، \babelsublr{NASA Johnson Space Center}، ملك عام؛ وخلية من \babelsublr{\emph{Emiliania huxleyi}}، أليسون ر.~تايلور (\babelsublr{Alison R.~Taylor})، \babelsublr{University of North Carolina Wilmington Microscopy Facility}، \babelsublr{CC~BY~2.5}؛ و\babelsublr{\emph{Orbulina universa}}، \babelsublr{H.~Takagi, K.~Kimoto, T.~Fujiki, H.~Saito, C.~Schmidt, M.~Kucera and K.~Moriya}، \babelsublr{CC~BY~4.0}؛ والدياتومات في مضيق مكمردو، غوردون ت.~تايلور (\babelsublr{Gordon T.~Taylor})، \babelsublr{Stony Brook University}، \babelsublr{NOAA Corps Collection}، ملك عام (كلها عبر \babelsublr{Wikimedia Commons}).
  \item الفصل~23 (الدياجينيز والأحواض الرسوبية): نموّات كوارتزية ثانوية في حجر بوتسدام الرملي، م.~ك.~رايغل (\babelsublr{M.~C.~Rygel})، \babelsublr{CC~BY-SA~4.0} (عبر \babelsublr{Wikimedia Commons}). سجل مستوى البحر المجمَّع عن \babelsublr{R.~M.~Spratt and L.~E.~Lisiecki} (2016)، \babelsublr{NOAA NCEI World Data Service for Paleoclimatology}؛ ومعاملات الانضغاط عن \babelsublr{J.~G.~Sclater and P.~A.~F.~Christie} (1980) كما وردت في \babelsublr{pyBacktrack} (\babelsublr{EarthByte}).
  \item الفصل~24 (السلّم الطبقي): المسمار الذهبي (\babelsublr{golden spike}) لـ\babelsublr{GSSP} الطابق السيلاندي في زوماييا، \babelsublr{Josugoni}، \babelsublr{CC~BY-SA~4.0}؛ وحدّ الكريتاسي والبليوجين في متنزه \babelsublr{Trinidad Lake State Park}، \babelsublr{Nationalparks}، \babelsublr{CC~BY-SA~2.5} (كلاهما عبر \babelsublr{Wikimedia Commons}). الأعمار القاعدية وألوان المخطط الكرونوستراتغرافي الدولي من \babelsublr{International Commission on Stratigraphy}، \babelsublr{CC~BY~4.0}.
  \item الفصل~25 (طرائق التأريخ الإشعاعي): بلورة زركون في التألّق الكاثودي، إيمانويل روكيت (\babelsublr{Emmanuel Roquette})، \babelsublr{CC~BY-SA~4.0}؛ ومطياف كتلة بالتأيّن الحراري، \babelsublr{Radiogenic}، \babelsublr{CC0} (كلاهما عبر \babelsublr{Wikimedia Commons}). منحنى المعايرة \babelsublr{IntCal20} عن \babelsublr{P.~J.~Reimer} وآخرين (2020)، \babelsublr{\emph{Radiocarbon}} 62، \babelsublr{CC~BY~4.0}، كما وُزّع مع حزمة \babelsublr{rintcal} (م.~بلاو، \babelsublr{M.~Blaauw}).
  \item الفصل~26 (بنية الغلاف الجوي والمحيط): استرجاع جهاز \babelsublr{CTD} و\babelsublr{rosette} في مضيق بيرينغ، \babelsublr{NASA Goddard Space Flight Center}، ملك عام (عبر \babelsublr{Wikimedia Commons}). درجة الحرارة والملوحة من \babelsublr{World Ocean Atlas 2023} (\babelsublr{NOAA National Centers for Environmental Information}).
  \item الفصل~27 (الإشعاع ونقل الحرارة): الأرض مضاءة بالشمس كما رُئيت من \babelsublr{Deep Space Climate Observatory}، \babelsublr{NASA Goddard Space Flight Center}، ملك عام (عبر \babelsublr{Wikimedia Commons}). الإشعاع والتدفقات السطحية والهطول من إعادة التحليل \babelsublr{NCEP-NCAR Reanalysis~1} (إ.~كالناي وآخرون، \babelsublr{E.~Kalnay and others}، 1996)، مُقدَّمة من \babelsublr{NOAA PSL}، \babelsublr{Boulder, Colorado, USA}، من موقعها على \babelsublr{https://psl.noaa.gov}.
  \item الفصل~28 (الدوران العام للغلاف الجوي): إعصار خارج المدارين في خليج ألاسكا، \babelsublr{NASA MODIS/LANCE}، ملك عام؛ ومنطقة التقارب بين المدارين من بيانات السحب \babelsublr{GOES-11}، \babelsublr{NASA Earth Observatory}، ملك عام (كلاهما عبر \babelsublr{Wikimedia Commons}). الرياح والضغط عند مستوى البحر من إعادة التحليل \babelsublr{NCEP-NCAR Reanalysis~1} (إ.~كالناي وآخرون، \babelsublr{E.~Kalnay and others}، 1996)، مُقدَّمة من \babelsublr{NOAA PSL}، \babelsublr{Boulder, Colorado, USA}، من موقعها على \babelsublr{https://psl.noaa.gov}.
  \item الفصل~29 (دوران المحيط): تيار الخليج كما يُرى من درجة حرارة سطح البحر، \babelsublr{NOAA NESDIS}، ملك عام؛ وعوّامة استكشاف عميق تُنزَل إلى الماء، صورة لمصوّر من \babelsublr{NOAA} لم يُذكر اسمه، ملك عام (كلاهما عبر \babelsublr{Wikimedia Commons}).
\end{itemize}
